\chapter{Supersymmetry Breaking and the Opening of Grand Unification}

These notes follow Leonard Susskind's ninth lecture in the supplementary course on supersymmetry and grand unification, with curation by LazyingArt LLC. The lecture announces its own route very clearly, and the chapter should follow that route just as clearly: first symmetry breaking in a familiar setting, then supersymmetry breaking in the simplest model from the previous lecture, then the effect of gravity, and finally the first opening move toward grand unification. The guiding idea is that symmetry breaking is first understood as the splitting of degenerate energy levels; only later do we reinterpret the same logic as the splitting of particle masses.

\section{Symmetry Breaking as Level Splitting}

Susskind begins from the most concrete possible place. Before speaking of supersymmetry at all, he wants a physical picture of symmetry breaking that one already trusts. He mentions the ferromagnet as the simplest model with the right qualitative structure: neighboring little magnetic moments prefer to align, and once they do, the system chooses a direction even though the underlying laws do not.

But before using the ferromagnet itself, he turns to an even more elementary example: an atom in a magnetic field. Suppose the atom has total angular momentum \(j\). Then there are
\begin{equation}
2j+1
\end{equation}
states, labelled by
\begin{equation}
m_J=-j,-j+1,\dots,j.
\end{equation}
In the absence of any preferred direction, rotational invariance forces these states to be degenerate:
\begin{equation}
E_{-j}=E_{-j+1}=\cdots=E_j.
\end{equation}
This is the first point the lecture wants to make. Symmetry shows up in the spectrum as degeneracy.

Now let a magnetic field pick out a direction. The rotational symmetry is broken by the environment, and the degeneracy is lifted. Susskind deliberately keeps the discussion schematic, but the pattern is the familiar Zeeman pattern,
\begin{equation}
\Delta E_{m_J}\propto B\,m_J.
\end{equation}
The exact coefficient is not the issue. What matters is that states which were identical in energy because a symmetry rotated them into one another are now split once the symmetry is broken. This is the template we will later read as boson--fermion mass splitting.

The lecture then makes the decisive move. Suppose the magnetic field is not imposed from outside, but is the field produced by a ferromagnet filling space. Then the atom again sees a distinguished direction and again sees its levels split. Yet this new direction is not built into the laws. The ferromagnet could have pointed anywhere. All uniformly aligned configurations have the same energy. The asymmetry is not in the equations; it is in the vacuum.

This is why the ferromagnet is such a powerful example. It simulates explicit symmetry breaking while actually being a case of spontaneous symmetry breaking. From the inside, the world looks asymmetric. From the outside, or more precisely from the point of view of the laws, it is not.

\subsection{Question \& Answer}

\paragraph{Question.} If one lives inside the broken phase, how can one tell the difference between spontaneous symmetry breaking and a world whose laws were asymmetric from the beginning?

\paragraph{Answer.} The answer is the existence of Goldstone modes. If we rotate the direction of magnetization everywhere by the same amount, that costs no energy at all. Now do something only slightly less global: let the direction of magnetization vary very slowly through space. The system does not remain frozen in that distorted configuration; it develops long-wavelength spin waves, or magnons.

The important point is that the energy of these modes becomes arbitrarily small when the wavelength becomes arbitrarily large:
\begin{equation}
E(k)\to 0 \qquad \text{as} \qquad k\to 0.
\end{equation}
Long wavelength means low momentum, and an excitation whose energy goes to zero as \(k\to 0\) is massless. Thus an observer inside the ferromagnet can perform experiments, study these spin waves, and infer that the underlying world is not explicitly asymmetric. The massless modes are the telltale sign of spontaneous symmetry breaking.

The lecture wants us to carry exactly this pattern forward. A broken vacuum is not merely a background with some numbers turned on. It comes with a diagnostic structure: degenerate states split, and broken generators acting on the vacuum create new low-energy states.

\section{Vacuum Criteria and Order Parameters}

Susskind next translates the ferromagnet into the language of the vacuum. Let \(M\) denote the magnetization field. If the system is so hot, or so fluctuating, that no direction is preferred on average, then
\begin{equation}
M=0,
\end{equation}
and the vacuum is symmetric. If any component of \(M\) is nonzero, then the vacuum has picked a direction:
\begin{equation}
M\neq 0.
\end{equation}
This is the order parameter. The laws remain rotationally invariant, but the vacuum does not.

He then gives the more formal statement. Let \(J\) denote the generator of rotations. Acting with \(J\) on a state gives the infinitesimal change produced by a small rotation. If the vacuum is symmetric, a small rotation changes nothing,
\begin{equation}
J|0\rangle=0.
\end{equation}
If the vacuum is not symmetric, then the same generator produces a genuinely new state,
\begin{equation}
J|0\rangle\neq 0.
\end{equation}
The lecture immediately interprets that state physically: it is the infinitely long wavelength Goldstone excitation,
\begin{equation}
J|0\rangle \sim |\text{Goldstone},\,k=0\rangle.
\end{equation}
This is a compact formal statement, but it carries the full physical content of the ferromagnet story. The vacuum is broken precisely when the symmetry generator does not annihilate it.

At this point the lecture has done all the preparatory work it needs. We now know what to look for when supersymmetry breaks: some component of the vacuum must play the role that \(M\) played for the ferromagnet, and the supersymmetry generator must fail to annihilate the vacuum. The only difference will be that supersymmetry generators are fermionic.

\section{Recap of the Chiral Superfield}

Only now does Susskind say, in effect: let us review supersymmetry. The point is not to introduce a new formalism from scratch. The point is to return to the one simple supersymmetric field theory already on the table and extract from it just what is needed for the vacuum analysis.

We begin with the chiral superfield \(\Phi\). In the form relevant here it depends effectively on \(y\) and \(\theta\), where
\begin{equation}
y^\mu=x^\mu+i\bar{\theta}\sigma^\mu\theta.
\end{equation}
Its component expansion is
\begin{equation}
\Phi(y,\theta)=\phi(y)+\psi(y)\theta+F(y)\theta^2.
\end{equation}
The field \(\phi\) is the scalar component, \(\psi\) the fermion component, and \(F\) the auxiliary field. Susskind reminds us that the chiral superfield must be complex; trying to make it real is not stable under supersymmetry.

\begin{figure}[h]
\centering
\includegraphics[width=0.62\textwidth]{lecture_09_figure_02.png}
\caption{Chiral superfield recap and the fermion kinetic term.}
\end{figure}

The basic superspace action from the previous lecture is
\begin{equation}
\int d^4\theta\,\Phi^\dagger\Phi.
\end{equation}
To integrate over the Grassmann variables means to pick out the highest \(\theta^2\bar{\theta}^2\) component of \(\Phi^\dagger\Phi\). Up to the numerical factors that Susskind intentionally suppresses, this produces the ordinary component Lagrangian
\begin{equation}
\mathcal{L}_{\mathrm{free}}
=
\partial_\mu\phi^\ast\,\partial^\mu\phi
+
\bar{\psi}\sigma^\mu\partial_\mu\psi
+
F^\ast F.
\end{equation}
The first term is the scalar kinetic term, the second is the usual fermion kinetic term, and the third is algebraic. There are no derivatives on \(F\). This is the structural fact on which the later derivation rests: \(F\) is not dynamical. It is present only to be eliminated.

\section{From the Superpotential to the Scalar Potential}

The lecture now adds the second ingredient. Last time the example was \(\Phi^2\), which generates a mass term. This time Susskind wants to go further, but characteristically he does so in stages. He first says: let us be a little more modest and take \(\Phi^n\), then later add different powers together into a general function.

There is a notational issue here. On the board in this stretch, the same Roman-looking letter is used too casually for two different things. To avoid confusion, we will write \(W(\Phi)\) for the superpotential and \(V(\phi)\) for the ordinary scalar potential that results after eliminating \(F\).

Start from
\begin{equation}
\int d^2\theta\,\Phi^n+\mathrm{h.c.}
\end{equation}
with
\begin{equation}
\Phi=\phi+\psi\theta+F\theta^2.
\end{equation}
What survives the \(d^2\theta\) integral? Only terms with exactly two \(\theta\)'s. There are only two kinds of such terms:

\begin{enumerate}
\item choose one \(F\theta^2\) and \(n-1\) factors of \(\phi\), giving \(nF\phi^{n-1}\);
\item choose two factors of \(\psi\theta\) and \(n-2\) factors of \(\phi\), giving \(\frac{n(n-1)}{2}\phi^{n-2}\psi\psi\).
\end{enumerate}

So one should read the board's arrow, quite literally, as
\begin{equation}
\int d^2\theta\,\Phi^n+\mathrm{h.c.}
\;\leadsto\;
nF\phi^{n-1}
+\frac{n(n-1)}{2}\phi^{n-2}\psi\psi
+\mathrm{h.c.}
\end{equation}
The lecture explicitly pauses to say that the arrow does \emph{not} mean equality. It means: these are the terms selected by the \(d^2\theta\) integral.

\begin{figure}[h]
\centering
\includegraphics[width=0.78\textwidth]{lecture_09_figure_03.png}
\caption{Extracting the two-\(\theta\) term from \(\Phi^n\).}
\end{figure}

Now the pattern becomes transparent. If \(W(\Phi)=\Phi^n\), then \(W'(\phi)=n\phi^{n-1}\). More generally,
\begin{equation}
\int d^2\theta\,W(\Phi)+\mathrm{h.c.}
\;\leadsto\;
F\,W'(\phi)
-\frac{1}{2}W''(\phi)\psi\psi
+\mathrm{h.c.}
\end{equation}
The lecture is not interested in the fermion bilinear at this moment, and neither are we. The vacuum question will be decided by the \(F\)-dependence.

\paragraph{Eliminating the auxiliary field.}
This is the moment at which the algebra tightens. Collect the terms containing \(F\). In cleaned notation they are
\begin{equation}
\mathcal{L}(F)
=
F\,W'(\phi)
+
F^\ast W'(\phi)^\ast
+
F^\ast F.
\end{equation}
The frame supports this structure, although the conjugation on the board is casual and locally unstable.

\begin{figure}[h]
\centering
\includegraphics[width=0.72\textwidth]{lecture_09_figure_04.png}
\caption{Auxiliary-field terms in the component Lagrangian.}
\end{figure}

Because \(F\) has no kinetic term, the Euler--Lagrange equation reduces to a purely algebraic condition:
\begin{align}
\frac{\partial \mathcal{L}}{\partial F}=0
&\quad\Rightarrow\quad
W'(\phi)+F^\ast=0,\\
\frac{\partial \mathcal{L}}{\partial F^\ast}=0
&\quad\Rightarrow\quad
W'(\phi)^\ast+F=0.
\end{align}
Thus
\begin{equation}
F^\ast=-W'(\phi),
\qquad
F=-W'(\phi)^\ast.
\end{equation}
Substituting back, all \(F\)-dependence collapses into a positive definite ordinary potential,
\begin{equation}
V(\phi)=|W'(\phi)|^2.
\end{equation}
This is the main mathematical payoff of the superspace calculation. Once the superpotential is given, the ordinary scalar potential is essentially read off by differentiation and squaring.

\paragraph{Worked examples.}
The lecture immediately uses this formula in the way it was meant to be used. If
\begin{equation}
W(\Phi)\sim m\Phi^2,
\end{equation}
then, ignoring the numerical factor of \(2\) that Susskind does not track carefully on the board,
\begin{equation}
W'(\phi)\sim m\phi,
\qquad
V(\phi)\sim m^2\,\phi^\ast\phi.
\end{equation}
This is just the ordinary scalar mass term.

If instead
\begin{equation}
W(\Phi)\sim g\Phi^3,
\end{equation}
then
\begin{equation}
W'(\phi)\sim g\phi^2,
\qquad
V(\phi)\sim g^2\,|\phi|^4.
\end{equation}
So a cubic superpotential produces a quartic scalar interaction. The lecture does not pursue the detailed Feynman rules here; it only wants the structural lesson: the ordinary potential follows directly from the superpotential.

\subsection{Question \& Answer}

\paragraph{Question.} How can a supersymmetric theory have a non-supersymmetric vacuum?

\paragraph{Answer.} Because the equations and the vacuum are not the same thing. The theory can be supersymmetric while its minimum is not.

The crucial fact is
\begin{equation}
V(\phi)=|W'(\phi)|^2\ge 0.
\end{equation}
Hence there are two qualitatively different possibilities.

If the minimum occurs at zero energy,
\begin{equation}
V_{\min}=0,
\end{equation}
then the quantity inside the absolute square must vanish:
\begin{equation}
W'(\phi)=0.
\end{equation}
Therefore
\begin{equation}
F=0.
\end{equation}
But \(F\) is the last component of the superfield. If it vanishes in the vacuum, then no superspace direction has been singled out, and the vacuum is supersymmetric.

If, on the other hand, the minimum is lifted above zero,
\begin{equation}
V_{\min}>0,
\end{equation}
then
\begin{equation}
W'(\phi)\neq 0,
\qquad
F\neq 0.
\end{equation}
This is the supersymmetric analogue of nonzero magnetization. The equations still possess the symmetry, but the vacuum has selected a direction in superspace. That is spontaneous supersymmetry breaking.

\begin{figure}[h]
\centering
\includegraphics[width=0.9\textwidth]{lecture_09_figure_05.png}
\caption{Vacuum condition and the positive supersymmetric potential.}
\end{figure}

The lecture's earlier review of the ferromagnet now pays off in full. A nonzero \(F\)-component plays the role of the order parameter. It is the thing in the vacuum that breaks the symmetry, and it is the thing that splits states that supersymmetry would otherwise keep together.

\section{Consequences of Supersymmetry Breaking}

Once we allow the vacuum to break supersymmetry, the lecture unfolds the consequences one by one rather than introducing more formalism.

The first consequence is vacuum energy. If \(V_{\min}>0\), the vacuum itself carries energy. In ordinary field theory one may ignore the additive normalization of the vacuum energy, but in a gravitational theory one cannot. The vacuum then gravitates. In modern language it acts as a cosmological constant. This is why Susskind suddenly asks, ``Who cares about vacuum energy?'' The answer is: gravity cares about it, and cares disastrously.

The second consequence is level splitting. The lecture deliberately returns to the magnetic analogy here. A nonzero \(F\)-term in the vacuum behaves, for the supermultiplet, the way the magnetic field behaved for the atom. It splits the energies, and therefore the masses, of the bosons and fermions that supersymmetry rotates into one another. The theory remains supersymmetric at the level of the equations; the solution does not.

The third consequence is the fermionic Goldstone mode. If the vacuum is not invariant, acting with the symmetry on the vacuum must produce a new zero-energy state. For ordinary symmetries this is a Goldstone boson. For supersymmetry it must be a fermion, because supersymmetry transformations change bosons into fermions. Thus broken supersymmetry implies a massless goldstino. Symbolically,
\begin{equation}
Q|0\rangle \sim |\text{goldstino},\,k=0\rangle.
\end{equation}
The lecture presents this not as a separate theorem but as the direct supersymmetric analogue of the earlier ferromagnet discussion.

\paragraph{Hidden sectors.}
Susskind then introduces the jargon by returning once more to the atom and the ferromagnet. If the magnetic field of the ferromagnet were enormous, the atom would be badly distorted and no longer recognizable. If the field is weak, the atom remains recognizable and its levels are only slightly split. That is the regime of interest for phenomenology. One sector does the symmetry breaking; another sector, weakly coupled to it, merely feels the small splittings. The symmetry-breaking sector is then called hidden.

This is not presented as a beautiful mechanism. Quite the opposite: the lecture emphasizes that realistic supersymmetry breaking tends to be contrived, involving extra fields and extra sectors. But the analogy is sharp. The hidden sector is the ferromagnet; the visible sector is the atom whose levels are split only slightly.

\paragraph{Gravity and supergravity.}
At this point the lecture becomes qualitative on purpose. Gravity changes the rigid supersymmetric potential. Susskind sketches an additional contribution and then immediately hesitates over the placement of the Planck mass. That hesitation matters: the chapter should not falsely canonize a formula the lecture itself treats as uncertain. The stable lesson is simply that coupling to supergravity modifies the vacuum-energy story and can shift the effective potential downward. In favorable circumstances this may rescue the theory from an unacceptably large vacuum energy.

\paragraph{The super-Higgs analogy.}
The final supersymmetry discussion comes back to a theme the lecture announced much earlier and then held in reserve: the relation between Goldstone modes and gauge fields. A massless spin-\(1\) particle has only two helicity states,
\begin{equation}
\lambda=\pm 1,
\end{equation}
whereas a massive spin-\(1\) particle has three,
\begin{equation}
\lambda=+1,0,-1.
\end{equation}
In the electroweak theory the missing longitudinal state is supplied by the Goldstone boson.

The lecture then applies the same logic to supergravity. A massless gravitino carries only the helicities
\begin{equation}
\lambda=\pm \frac{3}{2},
\end{equation}
but a massive spin-\(\tfrac{3}{2}\) particle must carry
\begin{equation}
\lambda=\pm \frac{3}{2},\ \pm \frac{1}{2}.
\end{equation}
The goldstino supplies the missing \(\pm \tfrac{1}{2}\) states. Thus the goldstino combines with the gravitino in exact analogy with the way a Goldstone boson combines with a gauge boson. This is the lecture's super-Higgs picture: the broken-supersymmetry Goldstone mode completes the degrees of freedom of a massive gravitino.

By this point Susskind declares the mathematics of supersymmetry finished. The lecture will return to its physics later, but the formal story needed for symmetry breaking has now been told.

\section{The First Step Toward \(SU(5)\)}

The pivot to grand unification is abrupt on purpose. The lecture says, in effect, that the hard abstract mathematics is over, and now we move to something easier: a little group theory. The change of subject should feel like a genuine new beginning.

The standard model gauge group is
\begin{equation}
SU(3)\times SU(2)\times U(1).
\end{equation}
The grand-unification question is whether this product group might be embedded in a larger, simpler group. The natural candidate is \(SU(5)\), precisely because it is the smallest simple group that contains \(SU(3)\times SU(2)\times U(1)\) as a subgroup.

By definition,
\begin{equation}
U^\dagger U=1,
\qquad
\det U=1
\end{equation}
for \(U\in SU(5)\). For an infinitesimal transformation,
\begin{equation}
U=1+i\epsilon G+O(\epsilon^2).
\end{equation}
Unitarity implies
\begin{equation}
G=G^\dagger,
\end{equation}
and the special condition implies
\begin{equation}
\operatorname{tr}G=0.
\end{equation}
So the generators of \(SU(5)\) are traceless Hermitian \(5\times 5\) matrices.

This immediately gives the counting rule. A Hermitian \(n\times n\) matrix carries \(n^2\) real parameters, and tracelessness removes one. Hence
\begin{equation}
\dim SU(n)=n^2-1,
\end{equation}
so in particular
\begin{equation}
\dim SU(5)=25-1=24.
\end{equation}
The transcript contains a local slip before this rule is stabilized, but the correct result is the one above.

The embedding itself is done by blocks. We split the \(5\times 5\) matrix as a \(2\oplus 3\) system. The three \(SU(2)\) generators \(\tau_a\) live in the upper-left \(2\times 2\) block, the eight \(SU(3)\) generators \(\lambda_A\) live in the lower-right \(3\times 3\) block, and one further diagonal generator commutes with both blocks. In the convention used in the lecture, one may write it as
\begin{equation}
Y\propto
\operatorname{diag}
\left(
-1,-1,\frac{2}{3},\frac{2}{3},\frac{2}{3}
\right).
\end{equation}
Its overall normalization is conventional; what matters here is its tracelessness and the different values it assigns to the \(2\)-block and \(3\)-block.

The lecture then makes its first particle assignment. The first two slots of the \(5\)-component column should be a weak doublet with no color, so they are naturally \(\nu_e\) and \(e^-\). The last three slots should carry color but no weak \(SU(2)\) quantum numbers. The tracelessness condition, read as a statement about charge sums inside a multiplet, then selects the anti-down quarks:
\begin{equation}
0+(-1)+\frac{1}{3}+\frac{1}{3}+\frac{1}{3}=0.
\end{equation}
Thus the first five-dimensional multiplet is
\begin{equation}
(\nu_e,\ e^-,\ \bar d,\ \bar d,\ \bar d).
\end{equation}
The three copies differ only by color.

And this is exactly where the lecture stops. The remaining ten particles of one family are not forced into the same representation. Instead Susskind says that they fit into a ten-dimensional representation of \(SU(5)\), and promises to return to that construction next time. The chapter should stop with the same incompleteness. Grand unification has not yet been developed; it has only been opened.

\section{Summary}

The lecture proceeds by analogy, but the analogy is never superficial. The Zeeman effect teaches us to see symmetry in the spectrum; the ferromagnet teaches us to distinguish the symmetry of the laws from the symmetry of the vacuum; the Goldstone mode teaches us how the broken vacuum announces itself. Only then do we return to the chiral superfield, derive
\begin{equation}
V(\phi)=|W'(\phi)|^2,
\end{equation}
and ask the central vacuum question. If the minimum lies at \(V=0\), then \(F=0\) and the vacuum is supersymmetric. If the minimum lies above zero, then \(F\neq 0\), the vacuum energy is nonzero, bosons and fermions split, and the goldstino appears. Gravity then turns vacuum energy from an innocuous shift into a physical problem, and supergravity offers a possible escape. Only after all of that does the lecture pivot to \(SU(5)\), where the first structural embedding of \(SU(3)\times SU(2)\times U(1)\) is set up and the first five-dimensional multiplet is identified.